\chapter{Атаки на ланцюг постачання програмного забезпечення}
\chaptershort{Атаки на ланцюг постачання}
\label{chap:ch15}

Класичні загрози (розділи~\ref{chap:ch01}, \ref{chap:ch05}) приходять ззовні, тому захист будують на межі довіри: недовірене блокують, довіреному дозволяють. \term{Атака на ланцюг постачання} підриває саму цю межу: ЗПК проникає \emph{через довірений канал} --- легітимне оновлення, підписаний інсталятор, популярну бібліотеку, --- яким захист довіряє за визначенням. Це робить такі атаки особливо небезпечними: вони обходять периметровий захист, уражають багато жертв одночасно (усіх користувачів скомпрометованого компонента) й зловживають найціннішим ресурсом безпеки --- довірою. Для України тема має особливу вагу: NotPetya, одна з найруйнівніших кібератак в історії, поширилася саме через скомпрометоване оновлення широко вживаного бухгалтерського ПЗ. Розділ розглядає поняття й анатомію атак на ланцюг постачання, їх типологію та кейси, механізми довіри й цілісності (підпис, походження, SLSA, SBOM, відтворювані збірки) та стратегії протидії, і завершується кейсом захисту конвеєра розробки. Тема поєднує моделювання загроз (розділ~\ref{chap:ch04}), розвідку (розділ~\ref{chap:ch12}) і реагування (розділ~\ref{chap:ch16}).

\section{Поняття й анатомія атаки}\label{sec:ch15-concept}

\term{Ланцюг постачання ПЗ} --- сукупність усіх компонентів, людей, процесів та інструментів, що беруть участь у створенні й доставці програмного продукту до користувача: вихідний код, сторонні бібліотеки, інструменти збірки, конвеєри CI/CD, репозиторії, механізми оновлення. Атака на ланцюг постачання~\cite{ch15:enisa} компрометує будь-яку ланку цього ланцюга так, щоб зловмисний код потрапив у кінцевий продукт і був доставлений жертвам через довірений канал.

\begin{tikzfig}
\begin{tikzpicture}[font=\scriptsize,
  b/.style={draw,rounded corners=2pt,minimum height=5.5mm,text width=15mm,align=center,font=\tiny},
  >={Stealth[length=1.4mm]}]
\node[b] (dev) at (0,0) {розробник / код};
\node[b] (dep) at (1.9,0) {залежності};
\node[b] (build) at (3.8,0) {збірка (CI/CD)};
\node[b] (sign) at (5.7,0) {підпис, реліз};
\node[b] (upd) at (7.6,0) {оновлення};
\node[b,fill=black!8] (vic) at (7.6,-1.3) {жертви};
\draw[->] (dev)--(dep); \draw[->] (dep)--(build); \draw[->] (build)--(sign); \draw[->] (sign)--(upd); \draw[->] (upd)--(vic);
\draw[->,dashed] (0,-1.1) node[font=\tiny]{компрометація} -- (dep.south);
\draw[->,dashed] (3.8,-1.1) -- (build.south);
\end{tikzpicture}
\figcap{Рис.~\thechapter.1 --- Ланцюг постачання ПЗ і точки можливої компрометації}
\end{tikzfig}

Анатомія атаки має спільну логіку: зловмисник компрометує точку з високим \emph{коефіцієнтом множення} (один скомпрометований компонент уражає тисячі жертв) і високою \emph{довірою} (жертви не перевіряють довірений канал). Саме поєднання множення й довіри робить ланцюг постачання стратегічно привабливою ціллю, особливо для APT (розділ~\ref{chap:ch13}): замість атакувати кожну жертву окремо, достатньо скомпрометувати спільного постачальника.

Коефіцієнт множення зручно виразити кількісно. Якщо скомпрометований компонент використовують $U$ прямих споживачів, а кожен, своєю чергою, постачає продукт далі в середньому $k$ наступним, то за $d$ рівнів поширення потенційна кількість уражених

\begin{equation}
V \approx U\sum_{j=0}^{d} k^{j} = U\,\frac{k^{d+1}-1}{k-1}\quad (k>1),
\label{eq:ch15-blast}
\end{equation}

зростає геометрично з глибиною ланцюга. Формула~\eqref{eq:ch15-blast} пояснює, чому компрометація популярної базової бібліотеки чи інструмента з великим $U$ і транзитивним поширенням ($k>1$) дає лавиноподібний «радіус ураження», непорівнянний із прямою атакою на окрему жертву, --- і чому саме вузли з найбільшим $U$ є пріоритетними як для зловмисника, так і для захисту.

\section{Типологія та кейси}\label{sec:ch15-typology}

Атаки класифікують за ураженою ланкою (табл.~\thechapter.1): компрометація \emph{оновлень} (підміна легітимного оновлення --- SolarWinds~\cite{ch15:solarwinds}, CCleaner~\cite{ch15:ccleaner}, NotPetya); компрометація \emph{залежностей} (зловмисний код у сторонніх бібліотеках відкритих репозиторіїв~\cite{ch15:ohm} --- typosquatting, захоплення покинутих пакетів, впровадження в популярні); компрометація \emph{інструментів збірки} й CI/CD (впровадження коду під час складання); компрометація \emph{підпису} (викрадення ключів для підписання зловмисного коду як легітимного).

\tabhead{Таблиця~\thechapter.1 --- Типологія атак на ланцюг постачання}

\begin{center}\footnotesize
\begin{tabularx}{\linewidth}{L{22mm} Y L{20mm}}
\toprule
Уражена ланка & Механізм & Приклад \\
\midrule
Оновлення & підміна релізу & SolarWinds, NotPetya \\
Залежності & зловмисний пакет & typosquatting \\
Збірка (CI/CD) & код під час складання & компрометація конвеєра \\
Підпис & викрадені ключі & підробка довіри \\
\bottomrule
\end{tabularx}
\end{center}

Кейс SolarWinds~\cite{ch15:solarwinds} --- взірець витонченості: APT впровадило бекдор у легітимне оновлення широко вживаного ПЗ мережевого керування, яке потрапило до тисяч організацій, зокрема урядових, через механізм оновлення з коректним підписом. NotPetya~\cite{ch15:notpetya-ua} --- взірець руйнівності: через оновлення бухгалтерського ПЗ код поширився лавиною й завдав мільярдних збитків світовій економіці. Ці кейси засвідчили: довірений підписаний канал не гарантує безпеки, якщо скомпрометоване саме джерело довіри.

Окремо варто виділити атаки на \emph{відкриті репозиторії залежностей}~\cite{ch15:ohm}, бо вони дешевші й масовіші за компрометацію великого вендора. \emph{Typosquatting} --- публікація зловмисного пакета з іменем, схожим на популярний (різниця в одну літеру), у розрахунку на помилку розробника при встановленні. \emph{Захоплення покинутих пакетів} --- перебирання контролю над занедбаним, але все ще вживаним пакетом через покинутий обліковий запис супровідника. \emph{Впровадження в легітимний пакет} --- компрометація облікового запису супровідника популярної бібліотеки й випуск отруєної версії. Спільна риса --- зловживання автоматизмом сучасної розробки: залежності підтягуються транзитивно, часто без ревізії, тож один отруєний пакет глибоко в дереві залежностей тихо потрапляє в тисячі продуктів. Це робить \emph{гігієну залежностей} --- фіксацію версій, перевірку походження, мінімізацію дерева --- базовою, а не факультативною практикою.

\section{Довіра, цілісність та походження}\label{sec:ch15-integrity}

Основою протидії є забезпечення \term{цілісності} й \term{походження} (provenance) артефактів --- гарантій того, що компонент не змінено й що він походить із довіреного процесу. Цифровий підпис гарантує цілісність і автентичність \emph{за умови} захищеності ключа; його недостатньо, якщо скомпрометовано збірку (підписується вже заражений артефакт) чи ключ. Тому сучасні підходи забезпечують цілісність \emph{усього} процесу: in-toto~\cite{ch15:intoto} фіксує й перевіряє кожен крок конвеєра; SLSA~\cite{ch15:slsa} визначає рівні зрілості гарантій цілісності складання; Sigstore~\cite{ch15:sigstore} спрощує підписання й прозорий облік артефактів; відтворювані збірки~\cite{ch15:reproducible} дозволяють незалежно перевірити, що бінарний артефакт відповідає вихідному коду.

Довіру до ланцюга зручно оцінювати як добуток довір до окремих ланок. Якщо ланцюг складається з $n$ ланок, кожна з надійністю (ймовірністю бути нескомпрометованою) $t_i$, то цілісність кінцевого артефакту за незалежних ланок

\begin{equation}
T_{\text{ланцюг}} = \prod_{i=1}^{n} t_i,
\label{eq:ch15-chain}
\end{equation}

не перевищує надійності найслабшої ланки й спадає зі зростанням довжини ланцюга. Формула~\eqref{eq:ch15-chain} формалізує інтуїцію «ланцюг міцний настільки, наскільки міцна його найслабша ланка» й пояснює, чому подовження ланцюга (більше залежностей, посередників, автоматичних кроків) системно знижує довіру, якщо кожну нову ланку не забезпечено гарантіями цілісності.

Гарантії цілісності складання прийнято впорядковувати за рівнями зрілості (табл.~\thechapter.2), що дозволяє поступово підвищувати надійність $t_i$ кожної ланки з~\eqref{eq:ch15-chain}.

\tabhead{Таблиця~\thechapter.2 --- Рівні гарантій цілісності складання (за моделлю SLSA)}

\begin{center}\footnotesize
\begin{tabularx}{\linewidth}{L{14mm} Y L{20mm}}
\toprule
Рівень & Гарантія & Механізм \\
\midrule
Базовий & документована збірка & опис процесу \\
Середній & походження артефакту & підписана provenance \\
Високий & ізольована, керована збірка & захищений конвеєр \\
Найвищий & незалежна перевірюваність & відтворювані збірки \\
\bottomrule
\end{tabularx}
\end{center}

\section{Прозорість: SBOM та інвентаризація}\label{sec:ch15-sbom}

Неможливо захистити те, про існування чого не відомо. Сучасне ПЗ містить десятки й сотні сторонніх компонентів, тому першою умовою захисту є \term{прозорість} складу. \term{SBOM} (Software Bill of Materials~\cite{ch15:sbom}) --- машиночитаний перелік усіх компонентів продукту та їх версій --- дозволяє швидко відповісти на критичне питання: «чи містить наш продукт уражений компонент?». Коли виявляють вразливість чи компрометацію популярної бібліотеки, організації зі SBOM миттєво визначають, що саме зачеплено, тоді як інші витрачають дні на ручний аудит.

\begin{tikzfig}
\begin{tikzpicture}[font=\scriptsize,
  b/.style={draw,rounded corners=2pt,minimum height=5.5mm,text width=18mm,align=center,font=\tiny},
  >={Stealth[length=1.4mm]}]
\node[b] (prod) at (0,0) {продукт};
\node[b] (sbom) at (2.6,0) {SBOM (склад)};
\node[b] (vuln) at (5.2,0) {база вразливостей / загроз};
\node[b] (act) at (7.8,0) {дія (патч, заміна)};
\draw[->] (prod)--(sbom); \draw[->] (sbom)--(vuln); \draw[->] (vuln)--(act);
\draw[->,dashed] (vuln) to[out=-90,in=-90] node[below,font=\tiny]{нова загроза} (sbom.south);
\end{tikzpicture}
\figcap{Рис.~\thechapter.2 --- SBOM як основа швидкого реагування на вразливості й компрометації компонентів}
\end{tikzfig}

SBOM --- не самоціль, а основа процесів: безперервного зіставлення складу з базами вразливостей і розвідкою загроз (розділ~\ref{chap:ch12}), керування оновленнями, оцінювання ризику залежностей. Він перетворює реагування на компрометацію компонента з хаотичного пошуку на кероване, швидке зіставлення, прямо скорочуючи час реагування~\eqref{eq:ch10-dwell}.

Повнота SBOM вимагає врахування \emph{транзитивних} залежностей, а не лише прямих. Сучасний продукт прямо залучає десятки бібліотек, але кожна з них тягне власні, ті --- свої, тож реальне дерево залежностей сягає сотень вузлів, більшість із яких розробник ніколи свідомо не обирав. Саме в цих глибоких, «невидимих» вузлах і ховаються отруєні пакети~\eqref{eq:ch15-blast}. Тому цінний SBOM рекурсивно розгортає все дерево й фіксує точні версії (а не діапазони), бо лише так зіставлення з розвідкою загроз дає однозначну відповідь, уражений продукт чи ні. Без транзитивної повноти SBOM створює оманливе відчуття контролю, покриваючи лише верхівку дерева.

\section{Стратегії протидії}\label{sec:ch15-defense}

Протидія атакам на ланцюг постачання багаторівнева. \emph{Зменшення й контроль залежностей}: мінімізація сторонніх компонентів, перевірка їх репутації й походження, фіксація версій, використання довірених дзеркал. \emph{Захист конвеєра збірки}: ізоляція, найменші привілеї, контроль цілісності (in-toto, SLSA), захист ключів підпису (апаратні модулі, короткоживучі ключі). \emph{Перевірка артефактів}: підпис і його верифікація, відтворювані збірки, прозорий облік. \emph{Прозорість}: SBOM, інвентаризація, моніторинг складу. \emph{Виявлення}: аномалії в поведінці оновлень і компонентів (навіть довірений компонент, що раптом звертається до нових серверів, підозрілий --- тут знову працюють поведінкові методи розділів~\ref{chap:ch07}, \ref{chap:ch10}).

Фундаментальний зсув --- від \emph{сліпої} довіри до \emph{перевірюваної}: замість «довіряти, бо підписано» --- «довіряти, бо перевірено цілісність усього процесу». Це втілення принципу нульової довіри (zero trust) у ланцюгу постачання: жодна ланка не є довіреною за замовчуванням, кожна підлягає перевірці. Повністю усунути ризик неможливо (сучасне ПЗ не будують без сторонніх компонентів), тому мета --- зробити компрометацію важчою, а наслідки --- виявлюваними й обмеженими. Українські дослідження систематизують загрози ланцюга постачання й методи забезпечення цілісності~\cite{ch15:supplychain-ua}.

Важливо усвідомлювати, що захист ланцюга постачання стикається з принциповим обмеженням виявлення: скомпрометований компонент \emph{за визначенням} проходить через довірений канал і часто має коректний підпис, тому сигнатурні й репутаційні методи (розділ~\ref{chap:ch06}) безсилі --- вони саме й призначені довіряти підписаному. Єдине, що лишається виявленню, --- поведінкові аномалії \emph{після} розгортання: довірений компонент, що раптом поводиться незвично (нові мережеві з'єднання, нетипові дії), викриває себе не походженням, а поведінкою. Це повертає до наскрізного принципу посібника: там, де статична довіра підведена, останнім рубежем лишається поведінковий моніторинг (розділи~\ref{chap:ch07}, \ref{chap:ch10}). Проте й він обмежений: витончена атака на кшталт SolarWinds навмисно діяла повільно й обережно, імітуючи легітимну активність, аби не викликати поведінкових сигналів, --- що знову підкреслює незамінність \emph{запобіжної} цілісності процесу над реактивним виявленням.

\section{Кейс: захист конвеєра розробки}\label{sec:ch15-case}

Розглянемо побудову захисту конвеєра розробки ПЗ від компрометації ланцюга постачання (захисна задача, за моделлю SLSA~\cite{ch15:slsa}).

\emph{Крок 1. Інвентаризація й SBOM.} Формуємо повний перелік компонентів і залежностей (SBOM~\cite{ch15:sbom}), встановлюємо безперервне зіставлення з базами вразливостей і розвідкою загроз.

\emph{Крок 2. Контроль залежностей.} Мінімізуємо й фіксуємо залежності, перевіряємо їх походження й репутацію, використовуємо довірені джерела, автоматично скануємо на відомі загрози й typosquatting.

\emph{Крок 3. Захист збірки.} Ізолюємо конвеєр CI/CD, застосовуємо найменші привілеї, фіксуємо й перевіряємо кожен крок (in-toto~\cite{ch15:intoto}), захищаємо ключі підпису апаратно й робимо їх короткоживучими.

\emph{Крок 4. Перевірка артефактів.} Підписуємо артефакти, забезпечуємо верифікацію підпису споживачами, за можливості впроваджуємо відтворювані збірки~\cite{ch15:reproducible} для незалежної перевірки відповідності коду.

\emph{Крок 5. Моніторинг і реагування.} Моніторимо поведінку власних оновлень і компонентів на аномалії, готуємо план швидкого реагування на компрометацію компонента (завдяки SBOM --- миттєве визначення ураженого) з повідомленням споживачів.

Кейс демонструє зсув від сліпої довіри до перевірюваної цілісності всього процесу: захист ланцюга постачання --- це не окремий засіб, а дисципліна цілісності, прозорості й перевірки на кожній ланці, доповнена поведінковим моніторингом і готовністю до реагування.

\section*{Висновки до розділу}
\addcontentsline{toc}{section}{Висновки до розділу}

Атака на ланцюг постачання підриває саму межу довіри, проводячи ЗПК через довірений канал --- оновлення, підписаний інсталятор, популярну бібліотеку, --- чим обходить периметровий захист і уражає багато жертв одночасно. Її сила --- у поєднанні множення (один компонент --- тисячі жертв) і довіри, що робить ланцюг стратегічною ціллю APT. Атаки класифікують за ураженою ланкою --- оновлення, залежності, збірка, підпис (табл.~\thechapter.1); кейси SolarWinds і NotPetya засвідчили, що підписаний канал не гарантує безпеки, якщо скомпрометоване джерело довіри. Протидія забезпечує цілісність і походження \emph{усього} процесу (in-toto, SLSA, Sigstore, відтворювані збірки), бо довіра ланцюга не перевищує надійності найслабшої ланки~\eqref{eq:ch15-chain}. Прозорість складу (SBOM) є основою швидкого реагування на компрометацію компонента. Стратегія протидії --- зсув від сліпої довіри до перевірюваної (нульова довіра в ланцюгу постачання): мінімізація й контроль залежностей, захист конвеєра й ключів, перевірка артефактів, прозорість і поведінковий моніторинг. Для України NotPetya є гірким прикладом руйнівного потенціалу таких атак. Головний висновок: повністю усунути ризик ланцюга постачання неможливо, але його знижують дисципліною цілісності, прозорості й перевірки на кожній ланці, роблячи компрометацію важчою, а наслідки --- виявлюваними й обмеженими.

\selfcheck
\begin{enumerate}
\item Чому атака на ланцюг постачання особливо небезпечна? Як вона підриває межу довіри?
\item Що таке ланцюг постачання ПЗ? Назвіть його ланки й точки можливої компрометації (рис.~\thechapter.1).
\item Поясніть роль множення й довіри в анатомії атаки. Чому ланцюг постачання привабливий для APT?
\item Опишіть типологію атак за ураженою ланкою (табл.~\thechapter.1).
\item Опишіть кейси SolarWinds і NotPetya. Який спільний урок вони дають щодо підписаного каналу?
\item Що таке цілісність і походження артефактів? Чому цифрового підпису недостатньо?
\item Які підходи забезпечують цілісність усього процесу (in-toto, SLSA, Sigstore, відтворювані збірки)?
\item Поясніть формулу довіри ланцюга~\eqref{eq:ch15-chain}. Чому подовження ланцюга знижує довіру?
\item Що таке SBOM? Яку критичну задачу він розв'язує при виявленні компрометації компонента?
\item Чому прозорість складу є першою умовою захисту? Як SBOM скорочує час реагування?
\item Назвіть рівні стратегії протидії атакам на ланцюг постачання.
\item Поясніть зсув від сліпої довіри до перевірюваної. Як принцип нульової довіри застосовується до ланцюга постачання?
\item Опишіть побудову захисту конвеєра розробки (кейс підрозділу~\ref{sec:ch15-case}). Чому захист --- це дисципліна, а не окремий засіб?
\item Дослідницьке: оцініть ризик залежностей обраного проєкту за SBOM, побудуйте граф залежностей і визначте найслабші ланки~\eqref{eq:ch15-chain}.
\item Дослідницьке: дослідіть ефективність відтворюваних збірок для виявлення впровадження коду на етапі складання на модельному конвеєрі.
\end{enumerate}
